\section{Policy Algebra}\label{sec:model}

\paragraph{Threat model.}
The LLM is policy-untrusted: it may propose arbitrary calls as a function of its visible transcript, whether from injected content~\cite{greshake2023not} or its own miscalibrated reasoning. The engine, engine configuration, registered authorities, and registered transformations constitute the trusted computing base; no theorem below depends on model behavior. We assume a durable, append-only event log, ensuring that policy checks evaluate against a consistent history. Addressing covert channels within model outputs is outside our scope~\cite{lampson1973note}.

\paragraph{Label state.}
The label state is $S = \powerset(U) \times T$, where $U$ is a finite set of opaque principal (reader) identities (decentralized reader sets~\cite{myers1997dlm}), $\powerset(U)$ its power set, and $T$ a finite trust chain. Audiences are specified as concrete static reader sets, with call placeholders resolved by direct substitution. Deployments can customize the trust chain $T$ for their environment. Because $S$ is a finite product lattice under component-wise meet ($\cap, \min$) and join ($\cup, \max$)~\cite{denning1976lattice,birkhoff1967lattice}, policy composition rules are associative, commutative, and monotonic by construction—guaranteeing deterministic state aggregation without custom proof obligations.

\paragraph{Checked actions.}
The checked monoid models policy actions on $S$ (cf.\ graded monads~\cite{katsumata2014graded}), equipped with action composition and identity. In the primary model, non-identity actions are strictly restrictive meets:
\[
  a_i = (\cap A_i, \min r_i).
\]
Rulings resolve execution policy requirements per call without modifying the label state directly. Because label aggregation ($\fold$, which accumulates sequential meet operations on an initial label) composes meets, it is commutative and monotone.

\begin{proposition}[Collapse]\label{prop:collapse}
In the primary model, folding an action sequence $a_1 \cdots a_n$ onto an initial label $L_0$ reduces to a single running meet:
\[
  \fold(a_1 \cdots a_n)(L_0) = L_0 \meet a_1 \meet \cdots \meet a_n,
\]
eliminating the need to re-evaluate historical trajectory events.
\end{proposition}
Proof in Appendix~\ref{app:proofs}.

\begin{proposition}[Monotone descent and settlement]\label{prop:descent}
In the primary model, every label update is a meet operation over a finite lattice. Consequently, trajectory labels descend monotonically under the partial order and stabilize after finitely many strict descents.
\end{proposition}
Note that label stabilization reflects bounded permission narrowing rather than execution termination, as identity actions may continue indefinitely. Proof in Appendix~\ref{app:proofs}.

\paragraph{Branch state.}
Execution trajectories form a directed \emph{trajectory tree} rooted at the primary parent trajectory. Labels are trajectory-local to each branch node, whereas the event log is globally shared. A \texttt{branch(p)} operation creates a child trajectory with initial label $L_c := L_p$ and a pre-declared exit contract. Subsequent child actions fold exclusively into $L_c$, leaving the parent label $L_p$ unchanged during child execution ($L_c \le L_p^{\mathrm{init}}$). Prior to an explicit \texttt{submit\_result(v)} merge, no child context or result values enter the parent trajectory.

\begin{definition}[Taint-confining branch]\label{def:taint-confining}
A branch is taint-confining when it is abandoned or exits under adequacy $L_p \le \lab(v)$, leaving the parent label $L_p$ preserved exactly (Theorem~\ref{thm:branch}, \S\ref{sec:branching}).
\end{definition}
Returning an un-sanitized restrictive raw value maintains safety but propagates restrictions back to the parent trajectory ($L_p' = L_p \meet \lab(v) \le L_p$), requiring re-evaluation against the pre-read acquisition policy.

\paragraph{The log monoid.}
The log is modeled as the free monoid over events under concatenation~\cite{pierce2002types}, where events are appended during execution lifecycle phases without separate approval phases. Events fall into two primary categories:
\begin{itemize}
  \item \textbf{Dispatch lifecycle facts} record audit details for tool execution attempts, including opened dispatches, failures ($\mathrm{CloseFailure}$), and indeterminate outcomes ($\mathrm{CloseIndeterminate}$).
  \item \textbf{Governance events} record rulings, dispatches, trajectory boundaries (turn ends, branch initializations, merges), narrowing acceptances, sanitizer applications, and label casts.
\end{itemize}

History predicates (\texttt{prior} and \texttt{no\_prior}) and Remedy planning evaluate over a narrower projection rather than raw facts: the \emph{committed-effect projection} $\mathcal{E}(\ell)$, a tree-wide projection over the shared event log incorporating a tool contract's \texttt{emits} sequence $K$. Effects commit exactly once during execution (either at a success checkpoint $\mathrm{DispatchSucceeded}(K)$ or a reported-success close $\mathrm{CloseSuccess}(K)$), where $\uplus$ denotes multiset union:
{\small
\[
  \mathcal{E}(\ell \cdot e) = \begin{cases}
    \mathcal{E}(\ell) \uplus K & \text{if } e \in \{\mathrm{DispatchSucceeded}(K), \mathrm{CloseSuccess}(K)\}, \\
    \mathcal{E}(\ell) & \text{otherwise (failure or indeterminate close).}
  \end{cases}
\]}
History predicates are defined directly over $\mathcal{E}$:
\[
  \mathrm{prior}(k, \ell) \iff k \in \mathcal{E}(\ell), \qquad \mathrm{no\_prior}(k, \ell) \iff k \notin \mathcal{E}(\ell).
\]
Table~\ref{tab:dispatch_outcomes} summarizes these outcome semantics. Under committed-effect semantics, $\mathrm{prior}(k)$ proves a committed effect $k$, while $\mathrm{no\_prior}(k)$ proves only the absence of a committed effect—not the absence of un-committed attempts.

\begin{table}[t]
\centering
\caption{Outcome semantics for the committed-effect projection $\mathcal{E}$.}
\label{tab:dispatch_outcomes}
\small
\renewcommand{\arraystretch}{1.15}
\begin{tabularx}{\columnwidth}{@{}l *{3}{>{\centering\arraybackslash}X} @{}}
\toprule
\textbf{Dispatch outcome} & \makecell[c]{\textbf{Committed}\\\textbf{to $\mathcal{E}$?}} & \makecell[c]{\textbf{Establishes}\\\textbf{$\mathrm{prior}(k)$?}} & \makecell[c]{\textbf{Invalidates}\\\textbf{$\mathrm{no\_prior}(k)$?}} \\
\midrule
Checkpoint / Success with $k$ & $\checkmark$ & $\checkmark$ & $\checkmark$ \\
Failure          & $\times$     & $\times$     & $\times$     \\
Indeterminate    & $\times$     & $\times$     & $\times$     \\
\bottomrule
\end{tabularx}
\end{table}

The engine queries the log solely to evaluate history predicates over $\mathcal{E}$, verify ruling validity, and perform audits. Cached projections, such as tracking observed event types, correspond to monoid homomorphisms~\cite{pierce2002types}. Separate ``effects dimensions'' used in prior architectures are represented directly as log projections rather than label components: an effect tracks execution history rather than data confidentiality and integrity. Unifying events and labels into a single monoid product would blur essential semantic distinctions between per-value policy bounds and per-run event logs.

\paragraph{Integrating external policy logic through log views.}
Custom execution policies and numerical thresholds are implemented via log projections coupled with external authority resolvers. For example, cumulative metrics such as financial spend limits (\texttt{finance.spend}) are computed by log views that determine whether an action requires human authorization. Recording these operations as named log events avoids introducing quantitative state into the label algebra, keeping scalar magnitudes outside the lattice model.

\paragraph{Contracts.}
A tool contract formalizes policy requirements and execution side effects for tool dispatches, specifying two output declarations and pre-dispatch preconditions:
\begin{itemize}
  \item \textbf{Label delta (\texttt{delta}):} An action declaring the Label contribution of the tool's raw output. Prospective evaluation verifies the expected label before execution, while trajectory folding applies the validated label of the value actually admitted (whether raw, sanitized, or cast-resolved).
  \item \textbf{Committed effect (\texttt{emits}):} A sequence of effect tokens $K$ appended to the log projection $\mathcal{E}$ upon reported dispatch success, whether or not the output value is admitted.
  \item \textbf{Preconditions (\texttt{requires}):} Policy preconditions:
  \begin{enumerate}
    \item \emph{Trust floor:} minimum trust score ($r \ge r_{\min}$);
    \item \emph{Audience cover:} $\readers(L) \supseteq \mathrm{recipients}$, ensuring current trajectory readers $L$ cover all intended output recipients (where $\readers(\cdot)$ extracts the reader set);
    \item \emph{Source bounds:} input context restrictions ($\subseteq C$);
    \item \emph{History predicates:} conditions over prior log events in $\mathcal{E}$ (negative \texttt{no\_prior} and positive \texttt{prior}); and
    \item \emph{Hard gates:} explicit authorization requirements.
  \end{enumerate}
\end{itemize}
Unlike history predicates, which are satisfied passively by past entries in $\mathcal{E}$, hard gates cannot be satisfied retroactively; they require an explicit ruling from an authorizing component during atomic execution (\S\ref{sec:rulings}). Routing \texttt{tags} convey no algebraic weight and serve solely to route requests to appropriate policy authorities.

Dynamic resolvers and variable placeholders belong to the trusted computing base. Surface conventions, such as deriving automatic source deltas from access control lists, simplify policy specification without altering underlying contract semantics.

\paragraph{Unknown.}
Unknown values are not represented as elements within the lattice carrier; there is no ordering such as $\mathit{trusted} < \mathit{unknown} < \mathit{suspicious}$. Unresolved attributes remain outside the label algebra, preserving total algebraic composition and monoid laws while policy checks detect missing information. Reclassifying unknown inputs requires explicit, audited cast events. A cast is defined either as a static constant or via a dynamic resolver restricted to a bounded set of admissible target labels. Static casts establish a fixed default policy, while target ceilings prevent dynamic classifiers from arbitrarily elevating un-trusted data. This treatment aligns with gradual typing interpretations of unknown values \cite{disney2011gradual,fennell2013gradual}, directly mirroring how static type checkers manage partially annotated Python or TypeScript codebases: unannotated inputs default to an unclassified \texttt{Unknown} boundary rather than an assumed default position in the type hierarchy, requiring explicit casts or assertions to enter the security algebra.

